\begin{frame}{Shift-chains}
    Given an \(r\)-uniform hypergraph \(H=(V,E)\) and a total ordering \(<_V\) on \(V\), each edge \(e\) can be interpreted as a vector:
    \(\vec{e}=\set{e_1,\dots,e_r}\) \\
    where \(e_1<e_2<\dots<e_r\).
    \begin{columns}[T]
    \begin{column}{0.52\textwidth}

      \medskip
      We define a partial order on edges by
      \vspace{-1em}
      \[\vec{e}<_E\vec{f}\Longleftrightarrow
        \forall_i\ e_i\leq f_i.\]

      \(\vec{e}\) and \(\vec{f}\) are comparable if either
      \vspace{-1em}
      \[\vec{e}<_E\vec{f} \text{ or } \vec{f}<_E\vec{e}\]
      
      A hypergraph is a shift-chain if \(<_E\) is a total order.
    \end{column}
    \begin{column}{0.4\textwidth}
      \centering
      \vspace{-3em}
      \only<1>{
        \begin{shiftchain}[
          column sep=5mm,
          row sep=4mm,
          dot radius=3.0pt,
          dot color=cfgr,
          tikz options={scale=1.1}
        ]{6}
          \shiftedge{1,2,3}
          \shiftedge{1,2,3}
          \shiftedge{1,3,4}
          \shiftedge{2,3,4}
          \shiftedge{3,4,5}
        \end{shiftchain}
      }
      \only<2->{
        \begin{shiftchain}[
          column sep=5mm,
          row sep=4mm,
          dot radius=3.0pt,
          dot color=cfblue,
          tikz options={scale=1.1}
        ]{6}
          \shiftedge{1,2,3,4}
          \shiftedge{1,2,3,5}
          \shiftedge{1,3,4,5}
          \shiftedge{2,3,4,6}
          \shiftedge{3,4,5,6}
        \end{shiftchain}
      }
      \vspace{3em}
      \onslide<3->{
        \begin{shiftchain}[
          column sep=5mm,
          row sep=4mm,
          dot radius=3.0pt,
          dot color=cfred,
          tikz options={scale=1.1}
        ]{7}
        \shiftedge{1,2,5,7}[$\vec{e}$]
        \shiftedge{3,4,5,6}[$\vec{f}$]
        \end{shiftchain}
        \par\smallskip
      \textcolor{cfred}{Not a shift-chain as}
      \par\vspace{-0.5em}
      \textcolor{cfred}{\(e_1<f_1 \land e_4 > f_4\)}
      }
    \end{column}
    \end{columns}
\end{frame}
